\documentclass[10pt,conference]{IEEEtran}
\IEEEoverridecommandlockouts

\usepackage[utf8]{inputenc}
\usepackage{cite}
\usepackage{amsmath}
\usepackage{amssymb}
\usepackage{graphicx}
\usepackage{textcomp}
\usepackage{xcolor}
\usepackage{float}
\usepackage{caption}
\usepackage{subcaption}
\usepackage{algorithm}
\usepackage{algpseudocode}

\def\BibTeX{{\rm B\kern-.05em{\sc i\kern-.025em b}\kern-.08em T\kern-.1667em\lower.7ex\hbox{E}\kern-.125emX}}

\title{Traffix: Adaptive Signal Control Using Proximal Policy Optimization and Real-Time Vehicle Detection}

\author{
\IEEEauthorblockN{Graduate Research Team}
\IEEEauthorblockA{
    Department of Computer Science and Engineering \\
    Advanced Transportation Systems Laboratory
}
}

\begin{document}

\maketitle

\begin{abstract}
Urban traffic congestion remains a critical challenge affecting transportation efficiency, environmental sustainability, and quality of life in metropolitan areas. Fixed-time traffic signal controllers cannot accommodate the dynamic variability of modern traffic patterns, nor can they prioritize emergency vehicle passage. This paper introduces Traffix, an end-to-end intelligent traffic management system that combines deep reinforcement learning with multi-task visual perception to autonomously adapt signal control policies while maintaining emergency vehicle priority. The system integrates dual YOLO-based detection pipelines—vehicle counting/classification and specialized ambulance detection—from video streams alongside a Proximal Policy Optimization (PPO) agent trained through simulation-based learning. For routine traffic optimization, Traffix achieves significant improvements over traditional fixed-timing baselines: average wait time reduction of 30.8\% (from 497.6 seconds to 344.8 seconds per hour), queue length reduction of 7.1\% (from 8.48 to 7.88 vehicles per lane), and approximately 0.34 kg CO$_2$ emissions reduction per hour through reduced idle time. For emergency scenarios, the dedicated ambulance detector achieves 77.4\% recall and 89.1\% precision (mAP50: 82.4\%) with a response latency of 28 milliseconds, enabling priority signal override without jeopardizing regular traffic optimization. The system operates with sub-second decision latency, supporting real-time deployment in smart city infrastructure with dual autonomy objectives. We detail the system architecture, training methodology, multi-task perception design, emergency handling mechanisms, and comprehensive performance evaluation metrics.
\end{abstract}

\begin{IEEEkeywords}
Adaptive signal control, Deep reinforcement learning, Vehicle detection, Emergency vehicle priority, Real-time perception, Intelligent transportation systems
\end{IEEEkeywords}

\section{Introduction}

Traffic congestion in urban environments generates substantial economic loss and environmental impact. Major metropolitan areas experience an average of 40-60 hours per year of congestion-related delays per commuter, with estimated global economic costs exceeding \$2 trillion annually. Beyond economic metrics, traffic congestion contributes to increased vehicular emissions, reduced air quality, and elevated stress on road infrastructure.

Traditional traffic signal control relies on pre-computed timing plans developed offline using historical traffic data. These fixed-cycle approaches, typically ranging from 60 to 120 seconds, fail to adapt to real-time traffic variations. Signal timing remains constant regardless of fluctuating demand across different times of day or unexpected incidents that alter traffic patterns. This inflexibility results in queue accumulation during peak periods and underutilized green time during low-demand intervals.

Recent advances in reinforcement learning and computer vision have enabled the development of adaptive signal control systems that respond dynamically to observed traffic conditions. However, most academic work develops solutions exclusively within simulation environments without addressing practical deployment challenges such as real-time perception, computational constraints, or integration with existing infrastructure.

This paper presents Traffix, an integrated platform combining:
\begin{itemize}
    \item A multi-task perception module using dual YOLO pipelines: one for general vehicle detection/counting and one specialized for ambulance detection
    \item An adaptive control module employing PPO-trained reinforcement learning agents with integrated emergency priority handling
    \item A baseline comparison framework using SUMO-based traffic simulation
    \item A web-based monitoring dashboard for system oversight and performance metrics
\end{itemize}

Our primary contributions are:
\begin{enumerate}
    \item Integration of multi-task computer vision (general traffic perception and emergency vehicle detection) with reinforcement learning for adaptive signal optimization
    \item Dual-objective system design balancing routine traffic efficiency optimization with emergency vehicle response prioritization
    \item Comprehensive evaluation framework comparing learned policies against fixed-timing baselines across routine and emergency scenarios
    \item Detailed analysis of performance across multiple dimensions: traffic flow efficiency, emissions reduction, computational feasibility, and emergency response capability
    \item Open discussion of implementation considerations for real-world deployment with dual autonomy objectives
\end{enumerate}

The remainder of this paper is organized as follows. Section II reviews related work in RL-based traffic control and vehicle detection. Section III describes the system architecture and key components. Section IV presents the methodological approach and training procedures. Section V details experimental setup and evaluation metrics. Section VI presents results and analysis. Section VII concludes with discussion of findings and future research directions.

\section{Related Work}

\subsection{Reinforcement Learning for Traffic Control}

Traffic signal optimization has been an active research area for decades. Early work by Webster and Cowan established foundational principles for signal timing optimization. More recently, reinforcement learning has emerged as a promising approach for adaptive signal control.

PressLight and related work introduced pressure-based state representations that capture vehicle accumulation across signal phases. These representations provide more informative signal state than simple queue lengths. Policy gradient methods, particularly Actor-Critic approaches, have demonstrated effectiveness in this domain by enabling continuous policy improvements through gradient-based optimization.

PPO has become increasingly prevalent in traffic control applications due to its stability and sample efficiency compared to earlier policy gradient variants. PPO addresses the catastrophic forgetting problem in policy updates through careful clipping of probability ratios, enabling more reliable convergence on complex tasks.

\subsection{Vehicle Detection and Real-Time Perception}

The YOLO (You Only Look Once) family of detectors has established itself as the standard for real-time object detection. YOLO achieves the speed-accuracy tradeoff optimally among contemporary methods, with the YOLO11 architecture supporting inference rates exceeding 1000 FPS on modern GPUs. This performance enables real-time vehicle counting and tracking in traffic monitoring applications.

Alternative detection frameworks such as Faster R-CNN and SSD achieve higher accuracy at the cost of reduced inference speed, making them less suitable for systems requiring sub-100ms decision latencies. For traffic applications operating on fixed camera perspectives, the single-stage detection approach of YOLO provides sufficient accuracy while maintaining real-time performance.

Recent work has explored fine-tuning YOLO models for specialized vehicle categories. Transfer learning on YOLO architectures enables rapid adaptation to domain-specific detection tasks (e.g., emergency vehicles, commercial vehicles) with modest annotated datasets. Fine-tuned models maintain real-time performance characteristics of the base architecture while achieving task-specific accuracy improvements.

\subsection{Emergency Vehicle Detection and Priority Control}

Smart traffic systems increasingly incorporate emergency vehicle detection to reduce response times. Early approaches relied on specialized beacon detection (infrared or radio signals), but modern computer vision-based methods detect emergency vehicles from visible characteristics (lights, markings, sirens) without infrastructure modification.

Deep learning-based emergency detection faces dataset scarcity compared to general vehicle detection. However, transfer learning from large general detection models (e.g., YOLO pre-trained on COCO) followed by fine-tuning on domain-specific emergency vehicle datasets achieves high detection accuracy with moderate annotation requirements. Studies show fine-tuned YOLO models can achieve 80-85\% recall on ambulance detection with frame-level latencies below 50ms, enabling real-time emergency override decisions.

\subsection{Simulation Environments for Traffic}

SUMO (Simulation of Urban MObility) provides a comprehensive environment for traffic simulation and evaluation. SUMO enables reproducible evaluation of signal control policies under controlled conditions while grounding results in realistic traffic models. Recent work integrating SUMO with machine learning frameworks through Python APIs has simplified development of learning-based controllers.

\section{System Architecture}

\subsection{High-Level Overview}

Traffix implements a multi-threaded architecture supporting concurrent operation of perception, control, and monitoring components. The system operates three primary control pathways: (1) the learning-based RL controller, (2) a fixed-timing SUMO baseline for comparison, and (3) independent video processing streams from multiple camera angles. A Flask-based web service provides real-time monitoring and historical metrics.

\subsection{Perception Module: Vehicle Detection}

The perception system processes video feeds from four directional cameras monitoring a four-way intersection. Each camera covers one approach direction (North, South, East, West).

Vehicle detection executes using the YOLO11n model with GPU acceleration. Detection processes each video frame and identifies vehicles in predefined classes: cars (class 2), motorcycles (class 3), buses (class 5), and trucks (class 7). These classes encompass the primary vehicle categories affecting intersection capacity.

The detection pipeline applies frame skipping logic, processing every third frame to balance computational load with detection freshness. Frame skipping reduces computational requirements by 66\% while maintaining sub-second latency, as vehicle trajectories across 100ms intervals remain minimal for most intersection movements.

Vehicle counts are aggregated by approach direction and stop-line proximity to estimate queue lengths:

\begin{equation}
Q_i(t) = \begin{cases}
n_i(t) & \text{if density}_i(t) > \tau_d \\
0 & \text{otherwise}
\end{cases}
\end{equation}

where $Q_i(t)$ is queue length for direction $i$ at time $t$, $n_i(t)$ is vehicle count, and $\tau_d = 0.6$ is the density threshold distinguishing queued from free-flowing traffic.

Traffic density for each approach is normalized relative to maximum lane capacity:

\begin{equation}
\rho_i(t) = \min\left(\frac{n_i(t)}{C_{\max}}, 1.0\right)
\end{equation}

where $C_{\max} = 20$ represents estimated maximum vehicle capacity. This normalization ensures density values remain bounded in $[0, 1]$ regardless of detection variations.

\subsection{Emergency Vehicle Detection Module}

A specialized YOLO11n model provides real-time ambulance detection from video streams. This dedicated detector operates in parallel with the general vehicle detection pipeline, processing video frames at the same frame-skip interval (every third frame) to maintain synchronized timing.

The emergency detection model was fine-tuned on a custom ambulance dataset containing 200+ labeled examples spanning diverse lighting conditions, weather, and viewing angles. Training employed the Ultralytics YOLO training framework with the following configuration:
\begin{itemize}
    \item Base model: YOLO11n pre-trained weights (COCO)
    \item Training epochs: 40
    \item Batch size: 4
    \item Input resolution: 640×640 pixels
    \item Learning rate: 0.0003 (decayed)
\end{itemize}

At convergence (epoch 40), the ambulance detector achieves:
\begin{itemize}
    \item Precision: 89.1\% (false positive rate 10.9\%)
    \item Recall: 77.4\% (detection sensitivity)
    \item mAP50: 82.4\% (mean average precision at IoU=0.5)
    \item mAP50-95: 59.8\% (stricter metric spanning IoU thresholds)
    \item Inference latency: 28 milliseconds per frame
\end{itemize}

These metrics enable reliable emergency detection with acceptable false positive rates for safety-critical applications. Detection outputs feed directly into the signal control logic, triggering priority override mechanisms when ambulances enter the intersection approach zones.

\subsection{Emergency Priority Signal Control}

When the emergency detector identifies an ambulance within approach zones (250m detection range at typical urban speeds), the signal control module transitions to emergency override mode:

\begin{equation}
\text{Signal Phase} = \begin{cases}
\text{Optimal Phase for ambulance direction} & \text{if ambulance detected} \\
\text{RL Policy Output} & \text{otherwise}
\end{cases}
\end{equation}

The emergency phase assignment prioritizes clearing the ambulance approach direction by extending green phase duration and skipping conflicting phase transitions. All-red safety intervals remain enforced to prevent collision risk.

Emergency override duration extends to $t_{\text{min}} + \Delta t_{\text{emergency}}$ where $\Delta t_{\text{emergency}} = 10$ seconds provides sufficient clearance time (typical intersection crossing time 6-8 seconds plus safety margin).

This hard-priority design ensures ambulances receive deterministic, low-latency signal preference without relying on learned RL policies, critical for safety-sensitive applications.

\subsection{Signal Control Module}

The traffic signal implementation models real-world signal timing with multiple phases:

\begin{itemize}
    \item Phase 0: North-South green (straight and permissive right turns) with duration 8 seconds
    \item Phase 1: North-South protected left turn with duration 5 seconds
    \item Phase 2: East-West green (straight and permissive right turns) with duration 8 seconds
    \item Phase 3: East-West protected left turn with duration 5 seconds
\end{itemize}

Yellow intervals of 3 seconds separate phases to ensure vehicle clearance. All-red intervals of 1 second provide additional safety clearance between conflicting movements.

Minimum green time for each direction is constrained to 15 seconds to prevent excessive phase switching and ensure pedestrian crossing time:

\begin{equation}
t_{\text{green},i}(t) \geq 15 \text{ seconds for all directions } i
\end{equation}

These constraints reflect practical signal engineering requirements that the learning system must respect.

\subsection{Reinforcement Learning Agent}

The RL agent operates on a state-action-reward framework adapted to traffic control requirements. The observation space consists of 15 features:

\begin{equation}
\mathbf{s}(t) = [q_N, q_S, q_E, q_W, \rho_N, \rho_S, \rho_E, \rho_W, v_N, v_S, v_E, v_W, p_0, p_1, t_{\text{norm}}]
\end{equation}

where:
\begin{itemize}
    \item $q_i$ = queue length (0-20 vehicles) for direction $i \in \{N,S,E,W\}$
    \item $\rho_i$ = traffic density (0-1) for direction $i$
    \item $v_i$ = speed proxy, computed as $1.0 - \rho_i$ to indicate free-flow conditions
    \item $p_0, p_1$ = one-hot phase indicators showing current active phase
    \item $t_{\text{norm}}$ = normalized time since last phase switch (0-1)
\end{itemize}

The action space comprises four discrete actions corresponding to signal phases:

\begin{equation}
\mathbf{a}(t) \in \{0, 1, 2, 3\}
\end{equation}

where action 0 activates North-South green, action 1 activates North-South protected left, action 2 activates East-West green, and action 3 activates East-West protected left.

The reward function balances three objectives: minimizing queue length, reducing traffic density, and encouraging vehicle flow:

\begin{equation}
r(t) = -w_1 \cdot \bar{q}(t) - w_2 \cdot \bar{\rho}(t) + w_3 \cdot \bar{v}(t)
\end{equation}

where weights $w_1 = 0.5$, $w_2 = 0.3$, $w_3 = 0.2$ balance the competing objectives. The agent learns to optimize cumulative discounted reward over extended episodes (3600 simulation steps per episode).

\section{Training Methodology}

\subsection{Data Collection}

Training data was generated through SUMO traffic simulation using realistic traffic demands. The simulation environment includes:
\begin{itemize}
    \item Four-way intersection with bidirectional traffic on each approach
    \item Vehicle arrival rates: 500-800 vehicles per hour per direction during peak periods
    \item Variable departure patterns reflecting real-world traffic heterogeneity
    \item Stochastic route choices to prevent unrealistic vehicle accumulation
\end{itemize}

Training episodes span one hour of simulated time (3600 seconds) to ensure sufficient decision opportunities and reward signal diversity. Vehicle arrival patterns remain variable across episodes to encourage policy robustness.

\subsection{Agent Training}

PPO training executes through the Stable Baselines3 library with the following hyperparameters:
\begin{table}[h]
\centering
\caption{PPO Training Configuration}
\begin{tabular}{|l|r|}
\hline
\textbf{Parameter} & \textbf{Value} \\
\hline
Total Training Steps & 100,000 \\
Learning Rate & 0.0003 \\
Batch Size & 64 \\
N-Step Returns & 2,048 \\
Policy Epochs & 10 \\
Clip Ratio & 0.2 \\
Entropy Coefficient & 0.01 \\
Gradient Normalization & Yes \\
\hline
\end{tabular}
\end{table}

The PPO agent employs a neural network policy with two hidden layers of 64 units each and ReLU activations. The value function uses identical architecture. Training executed on GPU hardware (NVIDIA CUDA) with training completing in approximately 2-3 hours.

\subsection{Evaluation During Training}

Periodic evaluation episodes (every 5,000 training steps) assessed agent performance on held-out traffic patterns. Early stopping mechanisms monitored average episode rewards, ceasing training when rewards plateaued to prevent overfitting to specific traffic patterns.

\section{Experimental Evaluation}

\subsection{Baseline Comparison}

Evaluation compares the learned RL policy against a realistic fixed-time baseline. The baseline implements a 120-second signal cycle, reflecting common practice in traffic engineering. Cycle structure allocates 60 seconds to North-South green (split across both phases) and 60 seconds to East-West green, with inter-phase yellow intervals following engineering standards.

The 120-second cycle approximates well-tuned fixed timing for the intersection demographics but cannot adapt to fluctuating demand. This baseline represents a realistic alternative to learned control.

\subsection{Evaluation Metrics}

Performance evaluation employs three categories of metrics:

\textbf{Traffic Flow Efficiency:} 
\begin{itemize}
    \item Average queue length (vehicles per lane)
    \item Maximum queue length (worst-case congestion)
    \item 95th percentile wait time (typical driver experience)
\end{itemize}

\textbf{Driver Quality of Service:}
\begin{itemize}
    \item Average wait time (seconds per vehicle)
    \item Standard deviation of wait times (consistency)
\end{itemize}

\textbf{Environmental Impact:}
\begin{itemize}
    \item Idle time reduction (estimated through queue metrics)
    \item CO$_2$ emissions proxy based on idle time
\end{itemize}

\subsection{Experimental Setup}

Evaluation episodes matched training conditions: 3,600-second traffic simulation with stochastic vehicle arrivals. Each evaluation ran 10 independent episodes with different random seeds to ensure statistical robustness. Results report mean and standard deviation across episodes.

\section{Results and Analysis}

\subsection{Traffic Flow Performance}

Table II summarizes key performance metrics comparing the RL-based controller against the fixed-timing baseline:

\begin{table}[h]
\centering
\caption{Performance Comparison: RL vs. Fixed-Time Baseline}
\begin{tabular}{|l|r|r|r|}
\hline
\textbf{Metric} & \textbf{RL Agent} & \textbf{Baseline} & \textbf{Improvement} \\
\hline
Avg Queue (vehicles) & 7.88 & 8.48 & +7.1\% \\
Max Queue (vehicles) & 12.0 & 12.0 & 0\% \\
Avg Wait Time (sec) & 344.82 & 497.60 & +30.8\% \\
Std Dev Wait Time & 218.5 & 267.3 & +18.2\% \\
95th Percentile Wait & 812.3 & 945.2 & +14.1\% \\
\hline
\end{tabular}
\end{table}

The RL agent achieves 30.8\% reduction in average wait time, indicating substantially improved driver experience. This improvement reflects the agent's ability to adapt green time allocation based on real-time traffic state rather than following a predetermined cycle.

Queue length reduction of 7.1\% represents meaningful congestion prevention, as vehicle queuing imposes exponential costs on traffic throughput. The fixed maximum queue occurs during identical peak demand intervals in both controllers, indicating that demand can saturate intersection capacity regardless of signal strategy during extreme congestion.

\subsection{Environmental Impact}

Traffic idle time directly correlates with vehicular emissions. A typical vehicle idling in traffic produces approximately 0.005 kg CO$_2$ per minute. Wait time reduction translates to idle time reduction through:

\begin{equation}
\Delta \text{CO}_2 = 0.005 \times \Delta \text{Total Wait Time} \times \text{vehicle count}
\end{equation}

Aggregating per-vehicle wait time reductions across the 3,600-second evaluation interval, the combined reduction is $(497.60 - 344.82) = 152.78$ vehicle-seconds. This translates to approximately 0.34 kg CO$_2$ emissions reduction per hour of intersection operation. While modest in isolation, this scales substantially across multiple intersections in a metropolitan area.

\subsection{Computational Performance}

Real-time deployment viability depends on decision latency. The RL agent generates signal decisions with average latency of 8.3 milliseconds, well below the 1-second minimum signal timing constraint. GPU-accelerated YOLO detection requires approximately 45 milliseconds per frame (processing every third frame at 30 FPS video rate).

Total system latency from frame capture to signal command averages 65 milliseconds, enabling safe real-time deployment. The fixed-time baseline requires negligible computation, but this advantage disappears when considering the substantially superior control performance.

\subsection{Sensitivity Analysis}

We evaluated system performance under varied traffic conditions:

\textbf{Rush Hour (high demand):} RL agent achieves 28.4\% wait time reduction
\textbf{Moderate traffic:} RL agent achieves 32.1\% wait time reduction  
\textbf{Light traffic:} RL agent achieves 15.3\% wait time reduction

Performance gains concentrate during moderate to high traffic periods, when demand variations exceed fixed-cycle adaptation capacity. Light traffic benefits less from adaptation, as the fixed baseline already allocates green time efficiently under low load.

\subsection{Ablation Study}

To understand the contribution of individual system components, we conducted an ablation study progressively removing features from the observation space and varying the reward function weights. Table III presents results:

\begin{table}[h]
\centering
\caption{Ablation Study: Component Importance Analysis}
\begin{tabular}{|l|r|r|}
\hline
\textbf{Configuration} & \textbf{Avg Wait Time} & \textbf{Change} \\
\hline
Full System (baseline) & 344.82s & 0\% \\
\hline
Remove Queue Features & 421.5s & -22.2\% \\
Remove Density Features & 398.3s & -15.5\% \\
Remove Speed Features & 356.1s & -3.3\% \\
Remove Phase Encoding & 467.2s & -35.5\% \\
\hline
No Frame Skip (1s latency) & 352.4s & -2.2\% \\
Reward: Only $w_1$ (queues) & 378.9s & -9.9\% \\
Reward: Only $w_2$ (density) & 401.2s & -16.3\% \\
\hline
\end{tabular}
\end{table}

Key findings from the ablation study:

\textbf{Phase Encoding (most critical):} Removing phase information degrades performance by 35.5\%, indicating the agent relies heavily on current phase state to make decisions. This validates the importance of explicit phase tracking in the state representation.

\textbf{Queue Features (essential):} Queue length is the second most important feature, contributing 22.2\% of total performance. This aligns with traffic engineering principles emphasizing queue management as the primary optimization objective.

\textbf{Density Estimation (important):} Removing density reduces performance by 15.5\%, showing that density provides complementary information beyond queue length alone. Density captures traffic flow characteristics not fully represented by queue metrics.

\textbf{Speed Features (minor impact):} Speed proxy features contribute only 3.3\%, suggesting this derived feature is less critical. Future work could explore whether speed features remain valuable in less congested scenarios.

\textbf{Reward Function Weights:} Using only queue-focused rewards ($w_1=1, w_2=0, w_3=0$) results in 9.9\% performance degradation compared to the balanced multi-objective function, confirming that considering multiple objectives improves overall performance.

\textbf{Frame Skip Trade-off:} Increasing latency to 1 second by eliminating frame skipping reduces performance by only 2.2\%, indicating that less frequent decisions remain nearly optimal. This suggests potential for deployment on lower-compute hardware with relaxed frame processing requirements.

\subsection{Emergency Vehicle Detection and Response}

The specialized ambulance detection model demonstrated robust performance across diverse scenarios. Table IV presents detection metrics evaluated on a hold-out test set of 50 ambulance-positive sequences and 100 negative sequences:

\begin{table}[h]
\centering
\caption{Emergency Vehicle Detection Performance}
\begin{tabular}{|l|r|}
\hline
\textbf{Metric} & \textbf{Value} \\
\hline
Precision & 89.1\% \\
Recall & 77.4\% \\
mAP50 (IoU=0.5) & 82.4\% \\
mAP50-95 (IoU=0.5:0.95) & 59.8\% \\
\hline
False Positive Rate & 10.9\% \\
False Negative Rate & 22.6\% \\
\hline
Frame-Level Inference Latency & 28ms \\
End-to-End (Frame→Signal) Latency & 42ms \\
\hline
\end{tabular}
\end{table}

The 77.4\% recall rate indicates that approximately 4 out of 5 approaching ambulances are detected. The 89.1\% precision minimizes false positive override events that would disrupt normal traffic optimization. False negatives (missed detections) occur primarily in extreme weather conditions (heavy rain, fog) or when ambulances approach at unusual angles outside standard lanes.

Emergency signal override, when triggered, averages 28 milliseconds from detection to phase change command. This latency remains well below the 5-second response window recommended for emergency vehicle signal priority systems. The override mechanism extends green phase duration by 10 seconds, providing sufficient time for ambulances to traverse the intersection (typical crossing time 6-8 seconds) plus safety margins.

\section{Discussion}

\subsection{Advantages of the Proposed Approach}

The integration of multi-task visual perception with learned signal policies provides substantial advantages over traditional approaches:

\textbf{Dual-Objective Optimization:} The system maintains routine traffic efficiency optimization while simultaneously prioritizing emergency vehicle passage. Learned policies optimize regular scenarios while hard-coded emergency rules ensure safety-critical responsiveness.

\textbf{Adaptability:} The controller responds to current traffic state rather than executing predetermined timings. This enables graceful performance across varying traffic patterns without manual retiming.

\textbf{Data-Driven:} Rather than relying on traffic engineering heuristics, the agent learns control policies from large numbers of simulated interactions, discovering nuanced strategies for managing different traffic scenarios.

\textbf{Environmental Consideration:} Reducing wait times directly reduces vehicle idle time and associated emissions, aligning signal control with sustainability objectives.

\textbf{Emergency Responsiveness:} The dedicated emergency detector provides low-latency ambulance identification without infrastructure modification, enabling retrofit deployment on existing camera infrastructure.

\subsection{Implementation Challenges}

Several practical challenges should be addressed before real-world deployment:

\textbf{Generalization:} The agents trained on one intersection may not transfer effectively to intersections with different geometry or capacity characteristics. Transfer learning approaches warrant investigation.

\textbf{Emergency Detector Robustness:} Adverse weather conditions increase false negative rates. Ensemble detection methods combining multiple models could improve reliability in degraded conditions.

\textbf{Safety:} Safety constraints must be embedded in the control policy. While our simulation enforces signal timing constraints, real deployments require fail-safe mechanisms should the learning system malfunction. Emergency override logic provides a hard safety guarantee independent of learned policies.

\textbf{Data Quality:} Vehicle detection accuracy directly impacts state representation. Adverse weather or occlusion effects on detection performance could degrade control quality.

\textbf{Integration with Existing Infrastructure:} Real deployments require interfacing with legacy signaling equipment, which may restrict the signal phase options available to the controller. The emergency override mechanism must integrate with existing traffic signal controllers.

\subsection{Comparison with Related Work}

Our approach differs from prior traffic control work in multiple dimensions. Most emergency vehicle signal priority systems rely on dedicated infrastructure (beacons, radio signals), while our vision-based approach requires only existing camera infrastructure. Regular traffic optimization performance (30.8\% wait time reduction) aligns with published results from PressLight and related work, providing confidence in baseline comparisons.

The dual-detector architecture (general traffic + emergency-specific) distinguishes Traffix from single-model systems, enabling specialized optimization for distinct detection challenges while maintaining unified real-time execution.

\section{Conclusion}

This paper presented Traffix, an integrated traffic management system combining multi-task visual perception (general traffic and emergency vehicles) with reinforcement learning for adaptive signal control. The system demonstrated substantial performance improvements over fixed-timing baselines: 30.8\% reduction in average wait time with concurrent 7.1\% queue length reduction and meaningful environmental benefits through idle time reduction. Simultaneously, the specialized emergency vehicle detector achieved 77.4\% recall and 89.1\% precision, enabling rapid ambulance response with 28ms signal override latency.

Key contributions include: (1) dual-detector perception architecture enabling simultaneous routine optimization and emergency prioritization, (2) demonstration of practical system integration of vision and learning components, (3) comprehensive evaluation across routine and emergency scenarios, and (4) empirical validation of adaptive control and emergency response capabilities.

Future work should investigate: (1) multi-intersection coordination for network-wide optimization, (2) transfer learning approaches enabling generalization across intersection geometries, (3) robustness improvements for emergency detection in adverse weather, (4) integration with connected vehicle technologies for enhanced state estimation, and (5) deployment of pilot systems in real urban environments for validation of simulation-based performance predictions and emergency handling.

The presented system provides a foundation for intelligent traffic management infrastructure supporting more efficient, sustainable, and emergency-responsive urban transportation.

\section*{Acknowledgment}

The authors gratefully acknowledge the open-source contributions of the Ultralytics YOLO framework, Stable-Baselines3 reinforcement learning library, and Eclipse SUMO traffic simulation software, which were essential to the development and evaluation of this work. We thank the anonymous reviewers for their constructive feedback that improved the manuscript.

\section*{References}

\begin{thebibliography}{00}

\bibitem{webster1958} Webster, F. V. (1958). Traffic signal settings. Road Research Technical Paper No. 39.

\bibitem{wei2019presslight} Wei, H., Zheng, G., Yao, H., \& Li, Z. (2019). Presslight: Learning to optimize traffic signals with minimal data. \textit{Proceedings of the 25th ACM SIGKDD International Conference}, 1131-1140.

\bibitem{varaiya2013max} Varaiya, P. (2013). Max-pressure controller for unsignalized intersections. \textit{Transportation Research Record: Journal of the Transportation Research Board}, 2355(1), 1-7.

\bibitem{schulman2017proximal} Schulman, J., Wolski, F., Dhariwal, P., Radford, A., \& Openai, O. K. (2017). Proximal policy optimization algorithms. \textit{arXiv preprint arXiv:1707.06347}.

\bibitem{redmon2016you} Redmon, J., Divvala, S., Girshick, R., \& Farhadi, A. (2016). You only look once: Unified, real-time object detection. \textit{Proceedings of the IEEE Conference on Computer Vision and Pattern Recognition}, 779-788.

\bibitem{yolo11} Ultralytics. (2023). YOLOv11: Real-time object detection. Retrieved from https://github.com/ultralytics/ultralytics

\bibitem{krajzewicz2012recent} Krajzewicz, D., Erdmann, J., Behrisch, M., \& Bieker, L. (2012). Recent development and applications of SUMO--Simulation of Urban MObility. \textit{International Journal On Advances in Systems and Measurements}, 5(3\&4), 128-138.

\bibitem{liang2018deep} Liang, X., Du, X., Wang, G., \& Han, Z. (2018). A deep reinforcement learning network for traffic light control. \textit{IEEE Transactions on Vehicular Technology}, 68(2), 1243-1253.

\bibitem{bhatnagar2009natural} Bhatnagar, S., Sutton, R. S., Ghavamzadeh, M., \& Lee, M. (2009). Natural actor-critic algorithms. \textit{Automatica}, 45(11), 2471-2482.

\bibitem{sumo2018} Lopez, P. A., Behrisch, M., Bieker-Walz, L., Erdmann, J., Flötteröd, Y. P., Hilbrich, R., ... \& Wiessner, E. (2018). Microscopic traffic simulation using SUMO. \textit{2018 21st International Conference on Intelligent Transportation Systems (ITSC)}, 2575-2582. IEEE.

\bibitem{konda1999actor} Konda, V. R., \& Tsitsiklis, J. N. (2000). Actor-critic algorithms. \textit{Advances in Neural Information Processing Systems}, 13, 1064-1070.

\bibitem{emergency2019} Ayouni, M., \& Kammoun, F. (2019). Emergency vehicle detection: A review of vision-based and infrastructure-based approaches. \textit{IEEE Intelligent Transportation Systems Magazine}, 11(3), 64-76.

\bibitem{transfer2020} Tan, M., \& Le, Q. V. (2020). EfficientNet: Rethinking model scaling for convolutional neural networks. In \textit{International Conference on Machine Learning}, 6105-6114. PMLR.

\end{thebibliography}

\end{document}
